The results of this work carry implications that extend beyond offline evaluation of a single model.
Time-series foundation models are increasingly deployed as components of \emph{agentic systems} — architectures in which a model perceives its environment, maintains beliefs about its state, explains its inferences to human operators, and triggers actions in response to detected conditions.
In such settings, a blind spot in the underlying forecaster is not merely an accuracy issue: it is a reliability risk that can propagate into flawed monitoring, incorrect explanations, and ultimately unsafe decisions.
This section discusses how the theoretical and empirical results of the preceding sections can inform the design of agentic systems that are robust to structural aliasing.

\subsection{Representing and Monitoring Aliasing Risks}

An agentic system operating over time-series data can maintain an explicit \emph{spectral risk model}: a structured representation that records, for each active signal stream, the proximity of dominant frequency components to the blind-spot set $\{m \cdot f_s / S\}_{m \geq 1}$.
This model can be derived analytically from the patch parameters of the underlying foundation model and updated dynamically as the spectral content of the incoming stream evolves.
When a signal enters a high-risk frequency band, the agent can flag its forecasts as potentially unreliable, trigger a switch to an alternative tokenisation configuration, or escalate to a human operator.

\subsection{Explaining Aliasing to Human Users}

Reliable human–agent interaction requires that the agent be able to communicate the \emph{reason} for its uncertainty, not merely the fact of it.
The ontological layer introduced in Section~\ref{sec:point3} provides the vocabulary for such explanations: rather than reporting a raw probing accuracy or a Bayes factor, the agent can state that ``the observed signal pattern is characteristic of [ontology class], which falls within a known blind-spot frequency band of the current tokeniser configuration, reducing confidence in the forecast by approximately [posterior HDI width] units.''
This framing connects the signal-theoretic analysis to domain-meaningful concepts, supporting informed human oversight.

\subsection{Bayesian Belief Updating and Decision Support}

\noindent\textbf{[To be completed.]}
This subsection will discuss how the Bayesian posteriors computed in Section~\ref{sec:point7} can be incorporated into the agent's belief state, enabling principled uncertainty propagation from the forecaster through to downstream decisions, and how the agent could trigger corrective actions — such as re-tokenisation or model switching — on the basis of posterior evidence thresholds.
